% Generated draft: requires mathematical and source review.
\documentclass[11pt]{article}
\usepackage[margin=25mm]{geometry}
\usepackage{fontspec}
\setmainfont{Latin Modern Roman}
\usepackage{amsmath,amssymb,mathtools,unicode-math}
\setmathfont{Latin Modern Math}
\usepackage{enumitem,needspace,xurl,hyperref}
\setlength{\parindent}{0pt}
\setlength{\parskip}{5pt}
\setlength{\emergencystretch}{3em}
\newcommand{\N}{\mathbb{N}}
\newcommand{\Z}{\mathbb{Z}}
\newcommand{\Q}{\mathbb{Q}}
\newcommand{\R}{\mathbb{R}}
\newcommand{\C}{\mathbb{C}}
\newcommand{\F}{\mathbb{F}}
\newcommand{\A}{\mathbb{A}}
\newcommand{\PP}{\mathbb{P}}
\newcommand{\E}{\mathbb{E}}
\newcommand{\eps}{\varepsilon}
\newcommand{\dd}{\,\mathrm{d}}
\begin{document}
\section*{UnsolvedMath ID: 30006000}
\textbf{Source problem code:} catalogue entry
\subsection{Definitions and mathematical statement}
The supplied record concerns the abelian Higgs energy, denoted by \(E_\epsilon\), and pairs \((u,\nabla)\), where \(u\) is a field and \(\nabla\) is a connection. A critical pair is a pair for which the first variation of the energy vanishes for every admissible variation. The record does not specify the energy formula, the underlying bundle, the admissible variations, or the regularity class.

The proposed assertion is the following schematic statement: for every integer \(n\geq 5\), every critical pair for \(E_\epsilon\) in the low-energy, nearly flat regime is gauge-equivalent to a two-dimensional solution. Here gauge equivalence refers to equivalence under the gauge transformations of the underlying abelian Higgs model. The record does not specify that model's gauge action. It also does not define the low-energy regime, the nearly flat regime, or the precise meaning of a two-dimensional solution. Thus this assertion cannot yet be expressed as a fully quantified mathematical conjecture. In particular, the domain, boundary or asymptotic conditions, parameter ranges, quantitative thresholds, and whether the conclusion is local or global remain unspecified.

The supplied original excerpt contains a separate, incomplete statement entitled ``Tilt-excess decay.'' Its hypotheses are recorded as follows: for every \(n\geq 3\) and sufficiently small \(0<\rho\leq\rho_0(n)\), there exist constants \(\epsilon_0(n,\rho)\) and \(\tau_0(n,\rho)\). The excerpt considers a critical pair for \(E_\epsilon\) on
\[
B_1^n=\{x\in\R^n:|x|<1\},
\]
subject to
\[
\epsilon\leq\epsilon_0(n,\rho),\qquad u(0)=0,\qquad
\int_{B_1^n}e_\epsilon(u,\nabla)\,\dd x
\leq 2\pi\omega_{n-2}+\tau_0(n,\rho).
\]
Here \(e_\epsilon\) is the energy-density notation appearing in the excerpt; its formula is not supplied. The conventional interpretation of \(\omega_{n-2}\) is the volume of the unit ball in \(\R^{n-2}\), but the excerpt does not explicitly define it. The excerpt ends at this hypothesis and supplies neither an excess functional nor a decay conclusion. These hypotheses therefore cannot be identified with a complete definition of the conjecture's low-energy, nearly flat regime.

The excerpt also states that ``Savin's result, as stated above'' holds for critical pairs in every dimension \(n\geq2\), while the reconstructed statement and supplied literature assessment leave general critical pairs in dimensions \(n\geq5\) open. The referenced preceding formulation of Savin's result is absent. Consequently, the relationship between that claimed theorem, tilt-excess decay, and the proposed two-dimensionality assertion is unresolved.
\subsection{Short English statement}
The proposed conjecture asks whether every low-energy, nearly flat critical abelian Higgs pair in dimension \(n\geq5\) becomes two-dimensional after a gauge transformation. The supplied record lacks the definitions and hypotheses needed for a precise formulation, and its original excerpt contains an incomplete theorem whose relationship to the conjecture is unclear. The supplied dated curation labels the problem partially solved; this draft does not independently verify that assessment.
\subsection{Sources}
\begin{itemize}
\item Supplied local UnsolvedMath registry record, permanent ID \(30006000\), title ``Two-Dimensionality of Low-Energy Abelian Higgs Fields,'' exact source problem code \texttt{catalogue entry}, in the collection ``Oberwolfach Reports Open Problems.'' The record lists proposed year 2024 and does not identify a proposer.
\item Supplied source citation: \emph{Calculus of Variations (2025)}, Oberwolfach Reports, DOI 10.4171/owr/2024/37, pp.\ 2113--2178. \url{https://doi.org/10.4171/owr/2024/37}. The citation's 2025 label and its 2024 DOI/report designation are retained as supplied.
\item Supplied literature source: ``Pigati--De Philippis--Halavati, abelian Higgs.'' \url{https://poisson.phc.dm.unipi.it/~pigati/papers/ymh-savin.pdf}. The record supplies no further bibliographic details for this reference.
\item Supplied literature source: ``Oberwolfach Report 37/2024.'' \url{https://ems.press/content/serial-article-files/50045}.
\item Literature assessment dated exactly 2026-08-22: the supplied curation states that low-energy two-dimensionality is proved for critical points in dimensions \(2\leq n\leq4\) and for local minimizers in every dimension, and remains open for general critical points in higher dimensions. This is supplied dated curation derived from the OWR material, not an independently verified literature assessment. The record expressly calls for independent verification of time-sensitive status and cited claims.
\item The registry labels the statement \texttt{corrected\_verified} and reports that it was reconstructed and checked against the official Oberwolfach report. These are supplied verification claims; no source was independently accessed for this draft. The excerpt attributes a viscosity approach to Savin and a variational proof of Savin's theorem to Wang, cited there as reference [5]; the bibliography identifying [5] is not supplied.
\end{itemize}
\end{document}
